\documentclass[../../main.tex]{subfiles}
\graphicspath{{\subfix{../../image/}}}

\begin{document}

\section{注释与参考文献}\label{section:ch08-8-参考文献}

\paragraph{\S8.1。}关于变分法的更多内容，见 Giaquinta [Gi] 与 Giaquinta--Hildebrandt [G-H]。变分方法的良好参考书为 Struwe [Str] 与 Zeidler [Zd, 第 3 卷]。

\paragraph{\S8.2。}\S8.2.2 中\cref{theorem:ch08-2-1-弱下半连续} 的证明取自 Ladyzhenskaya--Uraltseva [L-U]。\S8.2.5 基于 Ball--Murat（Pitman Research Notes Math.\ Series 254, Longman (1991), 25--37）。

\paragraph{\S8.3。}关于极小元的正则性（与部分正则性）的更多信息，见 Giaquinta [Gi]。

\paragraph{\S8.4。}Kinderlehrer--Stampacchia [K-S] 与 Friedman [Fr3] 的书解释了更多关于变分不等式与自由边界问题的内容。关于调和映射的更多信息，见 Hélein [Hl] 与 Lin--Wang [L-W]。

\paragraph{\S8.5。}山路定理归功于 Ambrosetti 与 Rabinowitz（J.\ Functional Analysis 14 (1973), 349--381）。关于鞍点方法的进一步结果，见 Rabinowitz [Ra]（\S8.5 的来源）。

\paragraph{\S8.6。}Noether 定理及其应用的讨论主要基于 Reed--Simon [R-S2] 与 Strauss [St1]。另见 Morawetz 对 Lax--Phillips [Lx-P] 的附录及 Lax [Lx3]。恒等式\eqref{align:ch08-6-eq24} 来自 Costa 与 Strauss（Quart.\ Appl.\ Math.\ (1981), 351--361）。

\paragraph{\S8.7。}习题 13：关于如何表述对偶变分原理，见 Ekeland--Temam [E-T]。习题 18 归功于 Alessandrini：见 Isakov [I, \S4.3]。

\end{document}
